\section{Grounding World Knowledge into Perception}
\label{sec:vlm-grounding}

To act in the physical world, an AI system must connect language-based world knowledge to perceptual observations. 
Vision-language models and multimodal large language models provide a first step toward such grounding by aligning language with images, videos, and visual scenes.

\subsection{Vision-Language and Multimodal Models}
VLMs and MLLMs extend LLMs with visual inputs, allowing them to perform visual question answering, image and video understanding, spatial reasoning, object localization, and instruction grounding. 
For Physical AI, these capabilities are important because agents must recognize objects, infer scene layouts, understand human instructions, and reason about affordances.

\subsection{Spatial, Temporal, and Affordance Grounding}
Physical AI requires more than object recognition. 
It requires understanding where objects are, how they relate spatially, what actions they afford, and how they may change over time. 
Multimodal grounding enables models to connect language concepts such as ``grasp,'' ``push,'' ``open,'' or ``avoid'' to visual evidence in the environment.

\subsection{The Language Bottleneck for Dense Physical States}
A major limitation of VLMs is that their outputs are often expressed in language. 
While language is effective for abstraction and communication, it is a lossy interface for dense physical states. 
Physical interaction requires representing trajectories, velocities, contact dynamics, occlusion, deformation, uncertainty, and long-horizon temporal evolution. 
This language bottleneck motivates models that can move beyond description toward action and prediction.